\section{Correct answers per kilowatt-hour}
\label{sec:accuracy}

Energy per token says nothing about whether an agent still gets its tasks right. Caching should not change
outputs in principle, but we did not assume runs with and without it are identical, so we measured accuracy directly on 400 graded multi-turn tasks of BFCL~v4~\cite{bfcl2025} (base and long-context), with the benchmark's own grader and Qwen3-32B on vLLM, and recorded the GPU energy of each whole run
(Table~\ref{tab:bfcl}).

\begin{table}[t]
\centering\small
\caption{BFCL v4 multi-turn, Qwen3-32B on vLLM, 400 tasks (200 base, 200 long-context), 100 at a time, one round. Correct: total (base\,/\,long-context). Per kWh: correct tasks per kWh of GPU energy.}
\label{tab:bfcl}
\begin{tabular}{lrrrr}
\toprule
 & Correct & Time & Energy & Per kWh \\
 & & (min) & (kWh) & \\
\midrule
GPU cache & 204 (113\,/\,91) & 64 & 0.619 & 330 \\
no cache & 201 (120\,/\,81) & 129 & 1.239 & 162 \\
\bottomrule
\end{tabular}
\end{table}

Accuracy differed by 3 of 400 tasks (204 against 201), with category-level differences of 7--10 tasks in opposite
directions; with one round we cannot separate a cache effect from run-to-run variation. With the cache the same graded work took half
the time and half the GPU energy, so it produced 2.0$\times$ the correct answers per kWh. We had predicted equal accuracy
within two points and at least 20\% less energy per correct task; the energy per correct task fell by 51\%. A second
round was dropped for budget reasons. In both runs 23--24 long-context tasks were rejected by the engine because the
conversation exceeded the 40,960-token context (HTTP 400) and were counted as incorrect, which lowers absolute accuracy for both policies alike.
